\documentclass[11pt,fontset=fandol]{ctexart}

\usepackage[a4paper,margin=1.15cm]{geometry}
\usepackage{array}
\usepackage{booktabs}
\usepackage{graphicx}
\usepackage{tabularx}
\usepackage{xcolor}

\setlength{\parindent}{0pt}
\renewcommand{\arraystretch}{1.18}
\pagestyle{empty}

\newcommand{\docTitle}{购买记录 PDF 示例}
\newcommand{\docDate}{2026年05月20日}

\newcommand{\typeRow}[4]{#1 & #2 & #3 & #4 \\}
\newcommand{\subtotalRow}[2]{\multicolumn{2}{l}{#1小计} & #2 & \\}

\newcommand{\topimage}[2]{
  \begin{minipage}[t][8.1cm][t]{0.485\linewidth}
    \centering
    \textbf{#1}\\[0.3em]
    \includegraphics[width=\linewidth,height=7.5cm,keepaspectratio]{#2}
  \end{minipage}
}

\newcommand{\invoiceimage}[2]{
  \begin{minipage}[t][14.6cm][t]{\linewidth}
    \centering
    \textbf{发票}\\[0.3em]
    \includegraphics[width=\linewidth,height=13.8cm,keepaspectratio,#1]{#2}
  \end{minipage}
}

\newcommand{\recordpage}[7]{
\clearpage
{\Large \textbf{#1}}\\[0.3em]
类型：#2 \qquad 金额：#3\\[0.7em]
\noindent
\topimage{购买记录}{#4}\hfill
\topimage{付款记录}{#5}

\vspace{0.7em}

\noindent
#6

\vfill
\textcolor{gray}{#7}
}

\begin{document}

\begin{center}
  {\zihao{2}\bfseries \docTitle}\\[0.35em]
  {\large \docDate}
\end{center}

\vspace{0.25em}

{\small
\begin{center}
\begin{tabularx}{0.99\linewidth}{@{}>{\raggedright\arraybackslash}p{1.8cm} X >{\centering\arraybackslash}p{2.1cm} >{\centering\arraybackslash}p{1.8cm}@{}}
\toprule
类型 & 物品名 & 价格 & 备注 \\
\midrule
\typeRow{道具}{得力文件夹}{234.36}{示例脱敏图}
\subtotalRow{道具}{234.36}
\midrule
\multicolumn{2}{l}{总计} & 234.36 & \\
\bottomrule
\end{tabularx}
\end{center}
}

\vspace{0.5em}

\textbf{说明：}本页展示内置构图。上方左侧为购买记录，右侧为付款记录，下方为发票。示例图片均已脱敏，金额保留用于演示核对逻辑。

\recordpage
{1. 得力文件夹}
{道具}
{¥234.36}
{../assets/examples/example_order_redacted.png}
{../assets/examples/example_payment_redacted.png}
{\invoiceimage{}{../assets/examples/example_invoice_redacted.png}}
{示例备注：订单金额、付款金额与发票价税合计一致，适合作为 skill README 中的演示样例。}

\clearpage
{\Large \textbf{TODO}}\\[0.8em]

\textbf{这个示例还可以继续补充}
\begin{itemize}
  \item 增加多类型条目，演示首页小计与总计的完整写法。
  \item 增加“缺发票”与“已有但未匹配”的占位示例。
  \item 如果改用 PDF 发票原件，可把发票页替换成 PDF 截图或分页插图。
\end{itemize}

\end{document}
